%%%%%%%%%%%%%%%%%%%%%%%%%%%%%%%%
\chapter{Nested Sweepout}\label{ch:nested-sweepout}

The construction of nested sweepouts is crucial for our work. As discussed in \cite[§6]{CL20}, nested families allow us to glue local deformations while keeping precise control over the maximal mass of the sweepout. A family $\{\Gamma_t\}_{t\in[0,1]}$ is \emph{nested} if it is given by the level sets of a $\partial$-transverse Morse function, so that the associated regions $\Omega_t = f^{-1}((-\infty,t))$ vary by elementary topological moves (handle attachments or deletions) as $t$ increases.

The main result of this chapter, Proposition~\ref{prop:nested-sweepout}, shows that any sweepout of a bounded relative open set $U$ can be approximated by a nested sweepout with arbitrarily small mass increase. The proof follows the strategy of \cite[Lemma~6.2]{CL20}, adapted to the free-boundary setting, and proceeds in three main steps:

\begin{enumerate}
\item \textbf{Step 1 (Preliminary Modification):} Subdivide the original sweepout $\{\Gamma_t\}_{t\in[0,1]}$
into a finite number of subintervals $[t_i, t_{i+1}]$ such that on each subinterval, the corresponding
regions $\Omega_t$ are either monotone increasing or monotone decreasing.

\item \textbf{Step 2 (Local Monotonization):} On each subinterval $[t_i, t_{i+1}]$, construct a nested
family connecting slightly modified versions of $\Omega_{t_i}$ and $\Omega_{t_{i+1}}$. This uses the
splitting and extension lemmas from Chapter~\ref{ch:splitting-and-extension} to build explicit Morse
functions whose level sets have controlled mass.

\item \textbf{Step 3 (Gluing):} Glue the nested families from adjacent intervals into a single global
nested sweepout. The construction ensures that adjacent families satisfy the compatibility condition
$\Omega_0^{i+1} \subset \Omega_1^i$, which allows for repeated application of a gluing proposition.
\end{enumerate}
\begin{proposition}\label{prop:nested-sweepout}
    Let $U \subset M$ be a bounded relative open subset with smooth relative boundary. Then for any $\eps > 0$, and any $\br{\Gamma_t} \in \mathcal{S}(U)$ with the corresponding open sets $\br{\Omega_t}$ with $\mathcal{F}(\br{\Gamma_t}) \leq A$, there exists a nested family $\br{\Tilde{\Gamma}_t}$ with the corresponding family of open sets $\br{\Tilde{\Omega}_t}$ such that $\Tilde{\Omega}_0 \subset \Omega_0$, $\Omega_1 \subset \Tilde{\Omega}_1$ with $\mathcal{F}(\br{\Tilde{\Gamma}_t}) \leq A + \eps$. In particular, this implies $W(U) = W_n(U)$.
\end{proposition}
The proof of this proposition follows the strategy of~\cite[Lemma 6.2]{CL20}, with the following modifications for the free-boundary setting:
\begin{enumerate}[label=\textup{(\roman*)}]
    \item \emph{Boundary balls}: When a ball $B_i$ in the covering meets $\partial M$, we work inside a fixed collar $\Phi: \partial M \times [0, \rho) \to M$ and regard $B_i$ as bilipschitz to a half-ball in $\mathbb{R}^{n+1}_+$. All local deformations are generated by vector fields in $\mathfrak{X}_{\mathrm{tan}}(M)$, preserving the free-boundary condition.
    \item \emph{Relative cycles}: The sweepout takes values in $\mathbf{Z}_n(M, \partial M;\mathbb{Z}_2)$ (see Remark~\ref{rmk:Z2}), and we use the relative coarea formula (Corollary~\ref{cor:coarea-filling}) to control boundary contributions.
    \item \emph{$\partial$-transverse Morse functions}: The Morse functions constructed in the preliminary modification are required to be $\partial$-transverse (Definition~\ref{def:boundary-transverse}), ensuring that level sets meet $\partial M$ transversally.
\end{enumerate}

\section{Step 1: Preliminary Modification}
Fix $\br{\Gamma_t} \in \mathcal{S}(U)$ with the corresponding open sets $\br{\Omega_t}$ with $\mathcal{F}(\br{\Gamma_t}) \leq A$. We apply the preliminary modification as in \cite[Lemma 6.2]{CL20} to $\br{\Gamma_t}$.
\begin{lemma}[Preliminary Modification]\label{lem:preliminary-modification}
    For any $\eps > 0$ there exists a partition $0 = t_0 < \dots < t_N = 1$ of $[0, 1]$ and a family $\br{\Gamma_t'}$ with the corresponding family of open sets $\br{\Omega_t'}$, such that the following holds:
    \begin{enumerate}
        \item[(1.1)] $\Omega_0' \subset \Omega_0$ and $\Omega_1 \subset \Omega_1'$;

        \item[(1.2)] $\mathcal{F}(\br{\Gamma_t'}) \leq \mathcal{F}(\br{\Gamma_t}) + \eps$;
        \item[(1.3)] For each $i = 0, \dots, N - 1$ we have one of the two possibilities:
        \begin{enumerate}
            \item[(A)] $\Omega_{t_i}' \subset \Omega_{t_{i + 1}}'$ and there exists a Morse function $g_i: \operatorname{cl}(\Omega_{t_{i + 1}}\backslash \Omega_{t_i}) \to [t_i, t_{i + 1}]$, such that $\Gamma_t' = g_i^{-1}(t)$ and $\Omega_t' = \Omega_{t_i}'\cup g_i^{-1}((-\infty, t))$ for $t \in [t_i, t_{i + 1}]$.
            \item[(B)] $\Omega_{t_{i + 1}}' \subset
             \Omega_{t_i}'$ and there exists a Morse function $g_i: \operatorname{cl}(\Omega_{t_i}\backslash\Omega_{t_{i + 1}}) \to [t_i, t_{i + 1}]$, such that $\Gamma_t' = g_i^{-1}(t)$ and $\Omega_t' = \Omega_{t_i}'\cup g_i^{-1}((-\infty, t])$ for $t \in [t_i, t_{i + 1}]$.
        \end{enumerate}
    \end{enumerate}
    \begin{proof}
        Let $M'$ be a compact subset of $M$ that contains the closure of $\Omega_t$ for all $t \in [0, 1]$.

        Set $L_{\mathrm{ball}}:=(1+\frac{\eps}{100W})^{1/n}-1$ so that $(1+L_{\mathrm{ball}})^n\le 1+\frac{\eps}{100W}$.

        Choose $r$ sufficiently small so that for every ball $B$ of radius less than or equal to $r$ in $M'$ the following holds:
        \begin{enumerate}
            \item[(i)] $B$ is $(1 + L_{\mathrm{ball}})$-bilipschitz diffeomorphic to the Euclidean ball of the same radius (or to a half-ball in $\mathbb{R}^{n+1}_+$ if $B$ meets $\partial M$);
            \item[(ii)] $\mathcal{H}^n(B\cap \Gamma_t) < \frac{\eps}{20}$ for $t \in [0, 1]$.
        \end{enumerate}
        Condition (ii) can be realized because of the no concentration of mass property~\ref{sw:no-concentration}.

        \noindent\emph{Boundary balls.}
        If a ball $B_i$ meets $\partial M$, we work inside a fixed collar
        $\Phi:\partial M\times[0,\rho)\to M$ and regard $B_i$ as bilipschitz to a half-ball
        in $\mathbb R^{n+1}_+$. All
        local deformations in such $B_i$ are generated by compactly
        supported vector fields in $\mathfrak{X}_{\mathrm{tan}}(M)$.
        This preserves the free-boundary condition and keeps the same
        bilipschitz constants as in (i).

        Let $\br{B_i}$ be a collection of $k$ balls of radius $r$ covering $M'$, such that balls of half the radius cover $M'$. We choose a partition $0 = s_0 < \dots < s_{N'} = 1$, such that
        \begin{enumerate}
            \item[(iii)] \begin{align*}
            &\mathcal{H}^{n + 1}(B_i\cap(\Omega_{s_j}\backslash\Omega_{s_{j + 1}})) + \mathcal{H}^{n + 1}(B_i\cap(\Omega_{s_{j + 1}}\backslash\Omega_{s_j})) \\
            &< \min\br{\frac{r\eps}{10k}, \left(\frac{\eps}{10}\right)^{\frac{n + 1}{n}}};
            \end{align*}
            \item[(iv)] $\abs{\mathcal{H}^n(B_i\cap \Gamma_{s_j}) - \mathcal{H}^n(B_i\cap \Gamma_{s_{j + 1}})} \leq \frac{\eps}{10k}$
        \end{enumerate}
        for each $j = 0, \dots, N'$ and $i = 1, \dots, k$.

        We define the new family $\br{\Gamma_t'}$ as follows. For $t = s_j$ we set $\Omega_t' = \Omega_t$ and $\Gamma_t' = \partial\Omega_t$ unless $\Gamma_t$ is a finite collection of points in which case we set $\Gamma_t' = \Gamma_t$ and $\Omega_t' = \varnothing$.

        Define a subdivision of $[s_j, s_{j + 1}]$ into $2k$ subintervals, $s_j = s_j^0 < \dots < s_j^{2k} = s_{j + 1}$. Let $\br{B_i'}$ be a collection of $k$ balls concentric with $B_i$ of radius between $r/2$ and $r$ and such that $\partial B_i'$ intersects $\Gamma_{s_j}$ and $\Gamma_{s_{j + 1}}$ transversally. Set
        \begin{equation*}
            U_j^1 = \Omega_{s_j}\backslash\Omega_{s_{j + 1}}\quad\text{and}\quad U_j^2 = \Omega_{s_{j + 1}}\backslash\Omega_{s_j}.
        \end{equation*}
        By the slicing theorem (Corollary~\ref{cor:coarea-filling}) and (iii), we may assume $B_i'$ satisfies
        \begin{equation}\label{eq:coarea-small-boundary-ball}
            \mathcal{H}^n(\partial B_i'\cap (U_j^1\cup U_j^2)) \leq \frac{\eps}{4k}.
        \end{equation}

        Since the balls of radius $r/2$ already cover $M'$, and each $B_i'$ has radius in $[r/2,r]$ and is concentric with $B_i$, the family $\{B_i'\}_{i=1}^k$ still covers $M'$. Inductively we define
        \begin{equation}\label{eq:omega-0-1-2}
            \begin{split}
                \Omega_{s_j^0}' &= \Omega_{s_j}'\\
                \Omega_{s_j^{2i - 1}}' &= \Omega_{s_j^{2i - 2}}'\backslash (B_i'\cap U_j^1)\\
                \Omega_{s_j^{2i}}' &= \Omega_{s_j^{2i - 1}}'\cup(B_i'\cap U_j^2)
            \end{split}
        \end{equation}
        for $i = 1, \dots, k$.

        Surfaces $\partial\Omega_{s_j^l}'$ may not be smooth, but there exists an arbitrarily small perturbation so that the boundaries are smooth. We perform these perturbations in the inward direction for $\Omega_{s_j^{2i - 1}}'$ and in the outward direction for $\Omega_{s_j^{2i}}'$. To simplify notation we do not rename the sets after the perturbations; since the perturbations are arbitrarily small all the estimates for masses and volumes remain valid.

        The following properties follow from the definition~\ref{eq:omega-0-1-2} and (i)-(ii):
        \begin{enumerate}
            \item[(a)] $\abs{\mathbf{M}(\partial_I\Omega_{s_j^l}') - \mathcal{H}^n(\Gamma_{s_j})} < \eps/2$;
            \item[(b)] $\Omega_{s_j^{2i - 1}}' \subset \Omega_{s_j^{2i}}'$ and $\Omega_{{s_j}^{2i - 1}}' \subset \Omega_{{s_j^{2i - 2}}}'$.
        \end{enumerate}
        We define $\Gamma_{s_j^l}' = \partial\Omega_{s_j^l}'$, unless $\Omega_{s_j^l}'$ is empty.

        \begin{figure}[h]
            \centering
            \includegraphics[width=0.8\textwidth]{figures/IMG_2670.jpg}
            \caption{Construction of $\Omega_{s_j^{2i - 2}}, \Omega_{s_j^{2i - 1}}, \Omega_{s_j^{2i}}$ for $i = 1$.}
            \label{fig:2670}
        \end{figure}

        If $\Omega_{s_j^l}'$ is empty we set $\Gamma_{s_j^l}'$ to be a point inside $\Omega_{s_j^{l - 1}}'$. By properties (iii) and (iv) surface $\Gamma_{s_j^l}'$ will satisfy the desired upper bound on the mass.

        To complete our construction we need to show the existence of two types of nested families:
        \begin{itemize}
            \item a nested family that starts on $\Gamma_{s_j^{2i - 1}}'$ and ends on $\Gamma_{s_j^{2i - 2}}'$;
            \item a nested family that starts on $\Gamma_{s_j^{2i - 1}}'$ and ends on $\Gamma_{s_j^{2i}}'$.
        \end{itemize}
        In both cases we want the homotopies to satisfy the desired upper bound on the masses.

        Consider the set $\Omega_{s_j^{2i - 2}}'\backslash\Omega_{s_j^{2i - 1}}' = B_i'\cap U_j^1$. After smoothing the corner we call this set $U$. We map $B_i'$ to $\R^{n + 1}$ by a $(1 + L_{\mathrm{ball}})$-bilipschitz diffeomorphism. Existence of the desired nested families follows by Lemma~\ref{lem:existence-nested-between-sets}. The upper bound for the mass follows from the upper bound at the endpoints and properties $(i)$ and $(ii)$.
        \begin{equation*}
            \begin{split}
                \mathcal{H}^n(\Gamma_t') \leq& \mathbf{M}(\partial_I\Omega_{s_j^{2i - 1}}') + \mathbf{M}(\partial_I(\Omega_{s_j^{2i - 2}}'\backslash\Omega_{s_j^{2i - 1}}'))\\
                &+ C(n)(1 + L_{\mathrm{ball}})^n\mathcal{H}^{n + 1}(\Omega_{s_j^{2i - 2}}'\backslash\Omega_{s_j^{2i - 1}}') + \eps\\
                \leq& \mathcal{H}^n(\Gamma_{s_j}) + \eps/2 + \eps/10 + \eps/4k\\
                &+ C(n)(1 + L_{\mathrm{ball}})^n\min\br{\frac{r\eps}{10k}, (\frac{\eps}{10})^{\frac{n + 1}{n}}}\\
                \leq& \mathcal{H}^n(\Gamma_{s_j}) + \eps
            \end{split}
        \end{equation*}
    \end{proof}
\end{lemma}
\section{Step 2: Local Monotonization}\label{sec:local-monotonization}

Assume that family $\br{\Gamma_t}$ satisfies conclusions of Lemma~\ref{lem:preliminary-modification} for the subdivision $0 = t_0 < \dots < t_N = 1$.

For every $\eps > 0$ and each $i = 0, \dots, N - 1$, we will inductively define sets $\Omega_0^i$ and $\Omega_1^i$ such that the following holds:
\begin{enumerate}
    \item[(2.1)] $\Omega_0^i \subset \Omega_1^i$;
    \item[(2.2)] $\max\br{\mathbf{M}(\partial_I\Omega_0^i), \mathbf{M}(\partial_I\Omega_1^i)}$ \\
    $\leq \max\br{\mathcal{H}^n(\Gamma_{t_i}), \mathcal{H}^n(\Gamma_{t_{i + 1}})} + \eps$;
    \item[(2.3)] $\Omega_{t_{i + 1}} \subset \Omega_1^i$ and $\Omega_0^i \subset \Omega_{t_i}$;
    \item[(2.4)] There exists a nested family of hypersurfaces $\br{\Gamma_t^i}$, $0 \leq t \leq 1$, with the corresponding family of nested open sets $\Omega_t^i$, such that
    \begin{equation*}
        \mathcal{H}^n(\Gamma_t^i) \leq \max\br{\mathcal{H}^n(\Gamma_{t_i}), \mathcal{H}^n(\Gamma_{t_{i+1}})} + \eps.
    \end{equation*}
\end{enumerate}
Assume (2.1) - (2.4) holds for all $\Omega_0^j$ and $\Omega_1^j$ for $j < i$. To define $\Omega_0^i$ and $\Omega_1^i$, we only need to consider the following two cases:
\begin{itemize}
    \item \textbf{Case (A) of Lemma~\ref{lem:preliminary-modification}~(1.3):}\\
    In this case, $\Omega_{t_i} \subset \Omega_{t_{i + 1}}$. We take $\Omega_0^i = \Omega_{t_i}$ and $\Omega_1^i = \Omega_{t_{i + 1}}$ so that (2.1)-(2.3) would be automatically satisfied. (2.4) would follow from case~(A) of Lemma~\ref{lem:preliminary-modification}~(1.3).
    \item \textbf{Case (B) of Lemma~\ref{lem:preliminary-modification}~(1.3):}\\
    In this case, $\Omega_{t_{i + 1}} \subset \Omega_{t_i}$. We define $\Omega_1^i = \Omega_{t_{i + 1}}$ and
    \begin{equation*}
        \Omega_0^i = \Omega_{t_{i + 1}}\backslash \operatorname{cl}(N_\delta(\partial\Omega_{t_{i + 1}}))
    \end{equation*}
    for sufficiently small $\delta > 0$ so that $N_\delta(\partial\Omega_{t_{i + 1}})$ is diffeomorphic to $\partial\Omega_{t_{i + 1}} \times [-\delta, \delta]$ and hypersurfaces equidistant from $\partial\Omega_{t_{i + 1}}$ in this neighbourhood and for each $s \in [-\delta, \delta]$, we have
    \begin{equation*}
        \mathcal{H}^n(\partial\Omega_{t_{i + 1}}\times \br{s}) \leq \mathbf{M}(\partial_I\Omega_{t_{i + 1}}) + \eps/2.
    \end{equation*}
    Thus, $\Omega_0^i \subset \Omega_{t_i}$ and we will set $\Omega_1^i = \Omega_{t_{i + 1}}$.
\end{itemize}
Consider $\Omega_{0}^{i + 1}$. In case (A), it is defined as $\Omega_{t_{i + 1}}$ thus
\begin{equation*}
    \Omega_{0}^{i + 1} = \Omega_{t_{i + 1}} = \Omega_1^i
\end{equation*}
In case (B), it is defined as
\begin{equation*}
    \Omega_{0}^{i + 1} \subset \Omega_{t_{i + 2}} \subset \Omega_{t_{i + 1}} = \Omega_1^i.
\end{equation*}
Therefore, we can conclude that
\begin{itemize}
    \item[(2.5)] $\Omega_0^{i + 1}\subset \Omega_1^i$.
\end{itemize}
Roughly speaking, the reason why (2.5) holds is because to construct $\Omega_0^{i + 1}$ we push $\Omega_{t_{i + 1}}$ inward (or not at all) and to construct $\Omega_1^{i}$ we push $\Omega_{t_{i + 1}}$ outwards (or not at all).
\section{Step 3: Gluing Nested Families}

The final step is to glue together the nested families constructed on each interval $[t_i, t_{i+1}]$
into a single global nested sweepout of $U$. The key observation from Step 2 is that adjacent
intervals satisfy the compatibility condition $\Omega_0^{i+1} \subset \Omega_1^i$, which allows us
to apply the gluing proposition below repeatedly.

In this section, we prove the following:
\begin{proposition}\label{prop:gluing-two-nested}
    Suppose $\br{\Gamma_t^a}$ and $\br{\Gamma_t^b}$ are two nested families (with corresponding families of open sets $\br{\Omega_t^a}$ and $\br{\Omega_t^b}$ respectively) and $\mathcal{H}^n(\Gamma_t^i) \leq W$. Suppose moreover that $\Omega_0^b \subset \Omega_1^a$. For any $\eps > 0$ there exists a nested family $\br{\Gamma_t}$ and a corresponding family of open sets $\br{\Omega_t}$ such that
    \begin{equation*}
        \mathcal{H}^n(\Gamma_t) \leq W + \eps,
    \end{equation*}
    $\Omega_1^b \subset \Omega_1$ and $\Omega_0 \subset \Omega_0^a$.
    \begin{proof}
        Let $\mathcal{S}$ denote the collection of all open sets $\Omega'$ such that $\Omega_0^b \subset \Omega' \subset \Omega_1^a$ and $\partial\Omega'$ is smooth. Let
        \begin{equation*}
            A = \inf_{\Omega'\in\mathcal{S}}\mathbf{M}(\partial_I\Omega')
        \end{equation*}
        and choose $\Omega\in \mathcal{S}$ with
        \begin{equation*}
            \mathbf{M}(\partial_I\Omega) < A + \frac{\eps}{4}.
        \end{equation*}
        We set $\alpha = \partial\Omega$. We claim that $\Omega$ and $\alpha$ satisfy properties (1) and (2) from Lemma~\ref{lem:extension}I for $\Omega_t = \Omega_t^a$. Indeed if $\Omega'$ satisfies $\Omega \subset \Omega' \subset \Omega_1$, we have
        \begin{equation*}
        \begin{split}
            \mathcal{H}^n(\alpha) <& A + \frac{\eps}{4}\\
            \leq& \mathbf{M}(\partial_I\Omega') + \frac{\eps}{4}
        \end{split}
        \end{equation*}
        Thus, we know that there exists a nested family $\br{\Tilde{\Gamma}_t^a}$ with the corresponding family of open sets $\br{\Tilde{\Omega}_t^a}$, such that $\Tilde{\Omega}_0^a \subset \Omega_0^a$, $\Tilde{\Gamma}_1^a = \alpha$ and
        \begin{equation*}
            \mathcal{H}^n(\Tilde{\Gamma}_t^a) \leq W + \eps.
        \end{equation*}
        On the other hand, $\Omega$ and $\alpha$ also satisfy properties $(1)'$ and $(2)'$ from Lemma~\ref{lem:extension}II for $\Omega_t = \Omega_t^b$. Indeed, for any open set $\Omega'$ with $\Omega_0^b \subset \Omega' \subset \Omega$ then again we have $\Omega_0^b \subset \Omega$ and
        \begin{equation*}
            \mathcal{H}^n(\alpha) < \mathbf{M}(\partial_I\Omega') + \frac{\eps}{4}
        \end{equation*}
        Therefore, there exists a nested family $\br{\Tilde{\Gamma}_t^b}$ with the corresponding family of open sets $\br{\Tilde{\Omega}_t^b}$, such that $\Omega_1^b \subset \Tilde{\Omega}_1^b$, $\Tilde{\Gamma}_0^b = \alpha$ and
        \begin{equation*}
            \mathcal{H}^n(\Tilde{\Gamma}_t^b) \leq W + \eps.
        \end{equation*}
        We define the desired nested family $\Gamma_t$ simply by concatenating these two nested families.
    \end{proof}
\end{proposition}
Now we complete the proof of Proposition~\ref{prop:nested-sweepout}. We apply local monotonization~\ref{sec:local-monotonization} to define families $\br{\Gamma_t^i}$ for $i = 1, \dots, N - 1$.

By (2.5), we have
\begin{equation*}
    \Omega_0^2 \subset \Omega_1^1
\end{equation*}
Hence, we can apply Proposition~\ref{prop:gluing-two-nested} to the nested family $\br{\Gamma_t^1}$ and $\br{\Gamma_t^2}$. We obtain a new nested family $\br{\Gamma_t^{1, 2}}$ with the corresponding family of open sets $\br{\Omega_t^{1, 2}}$. Moreover, by (2.5) and Proposition~\ref{prop:gluing-two-nested}, we have
\begin{equation*}
    \Omega_0^{1, 2} \subset \Omega_0^1 \subset \Omega_0 \quad\text{and}\quad\Omega_1 \subset \Omega_1^2 \subset \Omega_1^{1, 2}
\end{equation*}

Using (2.5) again, we have
\begin{equation*}
    \Omega_0^3 \subset \Omega_1^2 \subset \Omega_1^{1, 2}
\end{equation*}
Hence, we can apply Proposition~\ref{prop:gluing-two-nested} to the nested family $\br{\Gamma_t^3}$ and $\br{\Gamma_t^{1, 2}}$. We obtain a new nested family $\br{\Gamma_t^{1, 2, 3}}$ with the corresponding family of open sets $\br{\Omega_t^{1, 2, 3}}$.

We iterate this procedure inductively. At the $i$-th step, we can apply the Proposition~\ref{prop:gluing-two-nested} to families $\br{\Gamma_t^{1, \dots, i}}$ and $\br{\Gamma_t^{i+1}}$ to construct a new nested family $\br{\Gamma_t^{1, \dots, i, i + 1}}$ with $\Omega_0^{1, \dots, i}$ and $\br{\Omega_1} \subset \Omega_1^{1, \dots, i}$, which would again guarantee that $\Omega_0^{i + 2} \subset \Omega_1^{1, \dots, i}$, so that we can go to the next step.

After iteratively performing this operation $N$ times we obtain the desired nested family. This finishes the proof of Proposition~\ref{prop:nested-sweepout}.
